\documentclass[11pt]{article}
\usepackage[a4paper,margin=1in]{geometry}
\usepackage{amsmath,amssymb,amsthm}
\usepackage{graphicx}
\usepackage{xcolor}
\usepackage[hidelinks]{hyperref}
\usepackage{microtype}
\setlength{\emergencystretch}{2em}

\newtheorem{theorem}{Theorem}
\newtheorem{lemma}[theorem]{Lemma}
\newtheorem{corollary}[theorem]{Corollary}
\newcommand{\A}{\mathcal A}
\newcommand{\Hh}{\mathcal H}
\newcommand{\Kk}{\mathcal K}
\newcommand{\M}{\mathcal M}
\newcommand{\Tt}{\mathcal T}
\newcommand{\B}{\mathbb B}

\title{Homogeneous tails rule out strict adjoint cyclicity}
\author{Automated mathematical discovery run \texttt{fa\_banach\_001}}
\date{13 August 2026}

\begin{document}
\maketitle

\begin{center}
\fcolorbox{black}{gray!8}{\parbox{0.91\linewidth}{
\textbf{Status.} Candidate full solution, likely valid, subject to human
review.  For every regular unitarily invariant Hilbert function space in the
source, neither the full multiplier algebra nor any of its norm-closed
subalgebras has a strictly $*$-cyclic vector.  This settles the source's
characterization question and removes hypothesis (i) from its Theorem~5.11.
}}
\end{center}

\section{Source question}

The source is D.~P.~Blecher and R.~Clou\^atre, \emph{Null projections and
noncommutative function theory in operator algebras}, arXiv:2404.04788v2
\cite{Source}.  Let $\Hh$ be a regular unitarily invariant Hilbert space of
analytic functions on the ball $\B_d\subset\mathbb C^d$.  Write
$\M(\Hh)$ for its multiplier algebra and $\A(\Hh)$ for the multiplier-norm
closure of the polynomial multipliers.

A vector $\xi\in\Hh$ is \emph{strictly $*$-cyclic} for an operator algebra
$\mathfrak A$ if
\[
 \mathfrak A^*\xi=\{a^*\xi:a\in\mathfrak A\}=\Hh.
\]
This is exact equality, not merely density.  On source p.~24, immediately
after Theorem~5.11, the authors state that they do not know a characterization
of the spaces $\Hh$ for which $\A(\Hh)$ has no such vector.  They prove the
sufficient condition $\A(\Hh)\subsetneq\M(\Hh)$.  The source passage and the
downstream theorem are reproduced below.

\begin{center}
\includegraphics[width=\linewidth]{figures/open_problem_crop.png}
\end{center}

The answer within the stated class is uniform: all of these spaces have the
nonexistence property.  In fact, even $\M(\Hh)$ has no strictly $*$-cyclic
vector.

\section{Proof intuition}

Multipliers can only move Taylor mass upward in homogeneous degree.  Thus
every high-degree tail is invariant.  A fixed vector $\xi$ has arbitrarily
small norm in sufficiently high tails.

If $\mathfrak A^*\xi$ were all of $\Hh$, the open mapping theorem would give
a single constant $C$ such that every unit vector could be written as
$a^*\xi$ with $\|a\|\leq C$.  Apply this to a unit vector already in a very
high tail.  Pairing the equation $a^*\xi=y$ with $y$ forces the scalar $1$
to be bounded by $C$ times the tail norm of $\xi$, which tends to zero.  This
contradiction is independent of the size of the multiplier algebra and of
any commutant calculation.

\section{An invariant-filtration obstruction}

\begin{lemma}[Tail obstruction]\label{lem:tail}
Let $\mathfrak A\subset B(\Kk)$ be a norm-closed operator algebra.  Suppose
there are nonzero closed $\mathfrak A$-invariant subspaces
\[
 \Kk_1\supset\Kk_2\supset\cdots
\]
whose orthogonal projections $P_n$ converge strongly to zero.  Then no
$\xi\in\Kk$ is strictly $*$-cyclic for $\mathfrak A$.
\end{lemma}

\begin{proof}
Suppose $\mathfrak A^*\xi=\Kk$.  The bounded conjugate-linear map
\[
 Q_\xi:\mathfrak A\longrightarrow\Kk,
 \qquad Q_\xi(a)=a^*\xi,
\]
is surjective.  The open mapping theorem, applied equivalently to the
conjugate Banach space of $\mathfrak A$, gives $C>0$ such that every unit
vector $y\in\Kk$ has a preimage $a\in\mathfrak A$ satisfying
\begin{equation}\label{eq:uniform-preimage}
 a^*\xi=y,\qquad \|a\|\leq C.
\end{equation}

Choose a unit vector $y_n\in\Kk_n$ and an $a_n$ as in
\eqref{eq:uniform-preimage}.  Invariance gives $a_ny_n\in\Kk_n$, so
\[
 1=|\langle a_n^*\xi,y_n\rangle|
  =|\langle P_n\xi,a_ny_n\rangle|
  \leq C\|P_n\xi\|.
\]
But $P_n\xi\to0$, a contradiction.
\end{proof}

The strong convergence assumption can equivalently be replaced by
$\bigcap_n\Kk_n=\{0\}$ for a decreasing sequence of closed subspaces.

\section{Full answer for Hilbert function spaces}

For the source's $\Hh$, unitary invariance gives the orthogonal homogeneous
decomposition
\[
 \Hh=\bigoplus_{m=0}^{\infty}\Hh_m,
\]
where $\Hh_m$ consists of homogeneous polynomials of degree $m$.  Regularity
ensures that every $\Hh_m$ is nonzero.  Put
\[
 \Hh_{\geq n}=\bigoplus_{m\geq n}\Hh_m,
 \qquad P_{\geq n}=P_{\Hh_{\geq n}}.
\]
The projections $P_{\geq n}$ converge strongly to zero.

\begin{theorem}[No strict adjoint cyclicity]\label{thm:main}
Let $\Hh$ be a regular unitarily invariant Hilbert space of analytic
functions on $\B_d$.  If $\mathfrak A$ is any norm-closed subalgebra of
$\M(\Hh)$, then
\[
 \mathfrak A^*\xi\neq\Hh\qquad(\xi\in\Hh).
\]
In particular neither $\M(\Hh)$ nor $\A(\Hh)$ admits a strictly
$*$-cyclic vector.
\end{theorem}

\begin{proof}
Every multiplier preserves every $\Hh_{\geq n}$.  Indeed, if the Taylor
series of $f$ has no term below degree $n$ and $\varphi$ is analytic, then
the Taylor series of $\varphi f$ has no term below degree $n$.  Hence
\[
 M_\varphi\Hh_{\geq n}\subset\Hh_{\geq n}
 \qquad(M_\varphi\in\M(\Hh)).
\]
The same is therefore true for $\mathfrak A$.  Apply
Lemma~\ref{lem:tail} with $\Kk_n=\Hh_{\geq n}$.
\end{proof}

Thus the relation $\A(\Hh)\subsetneq\M(\Hh)$ used in the source is not
needed.  Even in a hypothetical equality case, the common algebra would
still have no strictly $*$-cyclic vector.

\section{Strengthening of the source theorem}

The source defines a Toeplitz $C^*$-algebra $\Tt(\Hh)$ generated by
$\A(\Hh)$ and a notion of $\M(\Hh)$-total nullity for projections in
$\Tt(\Hh)^{**}$.  Its Theorem~5.11 says that a closed projection $q$ with
$q\A(\Hh)=q\Tt(\Hh)$ is $\M(\Hh)$-totally null, hence
$\A(\Hh)$-null, provided that
\begin{enumerate}
\item $\A(\Hh)$ has no strictly $*$-cyclic vector; and
\item a nonempty $\M(\Hh)$-totally null set exists.
\end{enumerate}

\begin{corollary}[Theorem 5.11 without hypothesis (i)]\label{cor:null}
Let $\Hh$ be as in Theorem~\ref{thm:main}, and let
$q\in\Tt(\Hh)^{**}$ be a closed projection satisfying
\[
 q\A(\Hh)=q\Tt(\Hh).
\]
If there exists a nonempty $\M(\Hh)$-totally null set, then $q$ is
$\M(\Hh)$-totally null and hence $\A(\Hh)$-null.
\end{corollary}

\begin{proof}
Theorem~\ref{thm:main} verifies hypothesis (i) of source
Theorem~5.11 automatically.  Apply that theorem with its remaining
hypothesis (ii).
\end{proof}

No assertion is made when the source's existence hypothesis (ii) fails.

\section{Verification, limitations, and novelty audit}

Eight checks were made against the proof.  First, exact strict cyclicity was
kept distinct from dense/topological cyclicity.  Second, conjugate-linearity
causes no open-mapping issue after passing to the conjugate Banach space.
Third, high-degree tails are invariant under arbitrary analytic multipliers,
not merely polynomial norm limits.  Fourth, the homogeneous decomposition
is exhaustive, so its tail projections converge strongly to zero.  Fifth,
regularity supplies nonzero test vectors in arbitrarily high degrees.  Sixth,
the proof covers the commutant-equality edge that the source's sufficient
condition misses.  Seventh, no claim is made for truncated finite-dimensional
function spaces, where arbitrarily high nonzero tails are absent.  Eighth,
Corollary~\ref{cor:null} removes only source hypothesis (i) and preserves the
independent total-null-set hypothesis.

The result also follows directly from the quotient-map formulation: strict
$*$-cyclicity would make $Q_\xi$ a bounded surjection, whereas the unit
vectors $y_n\in\Hh_{\geq n}$ prove that its adjoint cannot be bounded below.
This supplies an independent functional-analytic reading of the same
contradiction.

A bounded novelty search used the exact source title, arXiv id, authors, and
the phrases ``strictly $*$-cyclic vector'', ``strictly star-cyclic'',
``multiplier algebra'', ``Hilbert function space'', and ``unitarily
invariant''.  The current arXiv v2 source, updated in September 2025, still
contains the open statement.  Located literature concerns forward strict
cyclicity of weighted shifts or general strictly cyclic algebras, not the
adjoint multiplier-algebra question.  No explicit later answer or this
homogeneous-tail argument was found.  Novelty confidence is therefore
provisional pending expert review.

\begin{thebibliography}{9}
\bibitem{Source}
D.~P.~Blecher and R.~Clou\^atre,
\emph{Null projections and noncommutative function theory in operator
algebras}, arXiv:2404.04788v2 (2025), especially Theorem~5.11 and p.~24.
\end{thebibliography}

\end{document}
